\documentclass[11pt]{article}

\usepackage[T1]{fontenc}
\usepackage{lmodern}
\usepackage[a4paper,margin=27mm]{geometry}
\usepackage{amsmath,amssymb,amsthm,mathtools,mathrsfs}
\usepackage{microtype}
\usepackage{enumitem}
\usepackage{booktabs,array,tabularx}
\usepackage{xcolor}
\usepackage[colorlinks=true,linkcolor=blue!48!black,
  citecolor=blue!48!black,urlcolor=blue!55!black]{hyperref}
\setlength{\emergencystretch}{2em}

\newtheorem{theorem}{Theorem}[section]
\newtheorem{proposition}[theorem]{Proposition}
\newtheorem{corollary}[theorem]{Corollary}
\newtheorem{lemma}[theorem]{Lemma}
\theoremstyle{definition}
\newtheorem{definition}[theorem]{Definition}
\theoremstyle{remark}
\newtheorem{remark}[theorem]{Remark}

\newcommand{\C}{\mathbb C}
\newcommand{\R}{\mathbb R}
\newcommand{\N}{\mathbb N}
\newcommand{\cA}{\mathcal A}
\newcommand{\cB}{\mathcal B}
\newcommand{\cG}{\mathcal G}
\newcommand{\cH}{\mathcal H}
\newcommand{\cI}{\mathcal I}
\newcommand{\cO}{\mathcal O}
\newcommand{\cQ}{\mathcal Q}
\newcommand{\cS}{\mathcal S}
\newcommand{\cW}{\mathcal W}
\newcommand{\id}{\operatorname{id}}
\newcommand{\im}{\operatorname{im}}
\newcommand{\Tr}{\operatorname{Tr}}
\newcommand{\dist}{\operatorname{dist}}
\newcommand{\Span}{\operatorname{span}}

\title{\textbf{The Polarized Conservativity Law}\\
\large Kernel-pair geometry, singular completion, and orientation at the Riemann divisor}
\author{Michael Haufschild}
\date{}

\begin{document}

\maketitle

\begin{abstract}
We prove a polarized conservativity law applying before any choice of
topology, Hilbert space, differential field, or cohomology theory.  For an
arbitrary source \(X\), observation \(\pi:X\to Q\), and target charge
\(\chi:X\to C\), the charge is recoverable from the attained observation
range exactly when
\[
       \pi(x)=\pi(y)\quad\Longrightarrow\quad\chi(x)=\chi(y).
\]
Serial postprocessing can only enlarge this kernel pair and can never repair
a lost charge; parallel observation repairs it exactly by intersecting
kernel pairs.  If \(\pi\) is reflection invariant and \(\chi\) is
reflection odd, factorization forces \(\chi=0\), while every sufficient
transverse test must separate each nonzero charged reflection pair.  Thus
all route-specific missing objects are realizations of one universal map: a
polarized factor \(d\) satisfying \(\chi=d\pi\).

For a split Green source the law becomes the unique observation decomposition
\[
 F=(F\jmath)T+(FD)H:
\]
absolute orientation plus relative Green transgression.  For a normed route,
continuous factorization becomes the domination inequality
\(\|\chi(x)\|\le C\|\pi(x)\|\).  This specialization has a sharp Hilbert
boundary: for injective dense-range maps \(P:G\to H_1\) and
\(Q:G\to H_2\), a proper closed bi-dense source hyperplane exists exactly
when neither map is bounded below.

Orientation itself determines the form of any repair.  If an involution
exchanges an idempotent with its complement, the grading \(2e-1\) is a
commutator.  Every cyclic scalar trace therefore annihilates it, whereas an
explicit graded two-dimensional trace detects it.  Hence no cyclic scalar
positive form can be the decisive Riemann observer.  A route that eventually
scalarizes must first retain orientation in a graded, noncyclic, or
vector-valued test; only completed routes additionally require uniform
separation in a resolved graph topology.

For the complete prime-normal Riemann source this boundary is quantitative.
Every finite-depth source map is injective and contains a divisor-faithful
channel, yet all fixed spectral points collapse to one cofinal source atom.
An informative closed countable-prime completion must therefore use a
topology strictly finer than finite-coordinate product observation.  If it
also sews a positive valuation sector to its reflected sector, its operator
norm grows beyond every finite power of the inverse Wold boundary volume.

The arithmetic Calder\'on--Witt source, based infinite algebraic jets, Gamma
canonical products, Bohr and ordered-pulse tomography, Stokes rigidity, and
cyclic orientation loss become native realizations of the same factorization
law.  The Riemann consequence is not another divisor equivalence: every
successful route must construct a source-derived orientation-complete factor;
completed prime-normal routes must make it nonlocal and beyond every finite
Rees order.  No zero-exclusion assertion is assumed or claimed.
\end{abstract}

\medskip
\noindent\textbf{Keywords.}
Riemann hypothesis; Green contraction; singular completion; orientation;
Gamma function; hypertranscendence; algebraic jets; cyclic trace.

\medskip
\noindent\textbf{2020 Mathematics Subject Classification.}
Primary 11M26, 47A06; Secondary 12H05, 14L15, 46C05.

\section{Introduction}

Many proposed routes to the Riemann hypothesis construct one of four things:
a source carrying the prime data, a transport into an analytic or Hilbert
space, a scalar positive form, or a divisor identity.  The persistent gap is
then described as positivity, self-adjointness, density, or a missing
functional equation.  Those descriptions are not equivalent.  They occur at
different stages, and passing between the stages can erase exactly the
orientation that locates a zero horizontally.

This paper identifies a common mechanism behind that loss.  Information is
not preserved merely because a transport is injective, has dense range, or
reconstructs every finite observable.  It is preserved after completion
only when the resolved test family is uniformly separating.  Before
completion, every observation of a split source divides canonically into an
absolute trace and a relative Green transgression.  After completion, a
proper closed sector can become dense in every visible leg exactly when all
those legs are singular.  Scalar cyclicization then creates a second,
independent loss: it annihilates commutator orientation even when the
underlying source remains perfectly resolved.

The point is not a new reformulation of the Riemann hypothesis.  The point is
a classification theorem for what a successful proof object must retain.
It answers four structural questions.

\begin{enumerate}[label=\textup{(\roman*)},leftmargin=2.2em]
\item What information exists at the source?  One absolute orientation line
      and a complete relative Green primitive.
\item When does completion erase source information?  Exactly when every
      chosen transport leg fails a lower bound.
\item Why do scalar traces and scalar positivity fail to orient the divisor?
      The orientation is a commutator and is therefore cyclically invisible.
\item What is the missing Riemann datum?  A source-derived,
      orientation-complete test in the native category of the route; if the
      route completes, this includes a uniform lower bound.
\end{enumerate}

These answers also separate genuine progress from decorative equivalences.
Within the present source framework, another regular source class is
unnecessary once the quotient is exact.  Another regular transform does not
repair the completed routes once phase and ordered-correlation tomography
are complete.  Another cyclic scalar positive quantity cannot detect a
commutator orientation.  A bounded Hilbert completion cannot repair a
two-sided singular defect.  Other native categories remain open, but all
must solve the same resolved-test problem.

Section~\ref{sec:law} states the route-independent polarized conservativity law.
Sections~\ref{sec:green}--\ref{sec:orientation} prove its algebraic,
topological, and graded parts.  Section~\ref{sec:source} identifies the
arithmetic and jet sources.  Section~\ref{sec:gamma} gives the Gamma
differential-algebraic shadow.  Section~\ref{sec:riemann} states the exact
Riemann research frontier and explains how several established positive and
negative results are consequences of the same law.

\section{The polarized conservativity law}
\label{sec:law}

We begin before choosing linearity, a Hilbert space, topology, differential
field, or cohomology theory.  Let \(X\) be a source, let \(\pi:X\to Q\) be
an observation, and let \(\chi:X\to C\) be the target charge.  The only
route-independent datum carried by \(\pi\) is its kernel pair
\[
       K_\pi=\{(x,y)\in X^2:\pi(x)=\pi(y)\}.
\tag{2.0}\label{eq:kernel-pair}
\]
We say that \(\pi\) is \emph{conservative for \(\chi\)} when
\(K_\pi\subseteq K_\chi\).

\begin{theorem}[Polarized conservativity law]
\label{thm:resolved-test}
For arbitrary sets and arbitrary observations:
\begin{enumerate}[label=\textup{(\roman*)},leftmargin=2.2em]
\item The charge \(\chi\) is recoverable from \(\pi\) on the attained
      observation range if and only if \(K_\pi\subseteq K_\chi\).
      Equivalently, there is a unique map
      \(d:\pi(X)\to C\) satisfying
      \[
                             \chi=d\pi.
      \tag{2.0a}\label{eq:polarized-factor}
      \]
\item A heterogeneous family \((\pi_j:X\to Q_j)_{j\in J}\) recovers
      \(\chi\) if and only if \(\chi\) factors through its joint record.
      The joint kernel pair is \(\bigcap_jK_{\pi_j}\).
\item Serial postprocessing can only enlarge a kernel pair.  If \(\chi\)
      does not factor through \(\pi\), it cannot factor through
      \(f\pi\) for any downstream map \(f\).  Injective postprocessing is
      neutral.  In contrast, adding a transverse channel \(\tau\) repairs
      recovery exactly when
      \[
                       K_\pi\cap K_\tau\subseteq K_\chi.
      \tag{2.0b}\label{eq:transverse-repair}
      \]
\item Let \(\sigma:X\to X\) be a reflection.  If \(\pi\sigma=\pi\) and
      \(\chi\sigma=-\chi\) with values in a characteristic-zero field,
      factorization through \(\pi\) forces \(\chi=0\).  Every sufficient
      transverse channel must separate \(x\) from \(\sigma x\) whenever
      \(\chi(x)\ne0\).
\item For linear maps \(E\xrightarrow{f}F\xrightarrow{g}G\) and a selected
      source sector \(S\subseteq E\), the composite is conservative on
      \(S\) exactly when \(f\) is conservative on \(S\) and \(g\) is
      conservative on \(f(S)\).  In normed categories, uniform lower bounds
      multiply along the chain.
\end{enumerate}
\end{theorem}

\begin{proof}
If \(\chi=d\pi\), equality of observations implies equality of charges.
Conversely, if \(K_\pi\subseteq K_\chi\), define
\(d(\pi x)=\chi(x)\); kernel-pair refinement makes this definition
independent of the representative, and uniqueness on \(\pi(X)\) is
immediate.  This proves (i).  Apply (i) to the joint record to obtain (ii).
For (iii), \(\pi(x)=\pi(y)\) implies
\(f\pi(x)=f\pi(y)\), whereas the joint record has kernel pair
\(K_\pi\cap K_\tau\).  If \(\pi\) is reflection invariant, then
\(\chi(\sigma x)=\chi(x)\) under factorization, while oddness gives
\(\chi(\sigma x)=-\chi(x)\); characteristic zero yields \(\chi(x)=0\).
A repair must therefore separate every nonzero charged reflection pair.
Finally,
\(\ker(gf)=f^{-1}(\ker g)\).  Intersecting with \(S\) gives the two
conservativity conditions in (v); multiplying the two lower-bound
inequalities gives the quantitative statement.
\end{proof}

The map \(d\) is the universal missing object.  Depending on the route, it
appears as a Green retraction, Stokes reporter, differential-field
reconstruction, graded trace, positive contraction, or unbounded inverse.
Theorem~\ref{thm:resolved-test}(iii) is the irreversibility law: once a stage
identifies a pair carrying different orientation charge, no downstream
consumer can restore it.

For linear tests on a resolved quotient, the theorem reduces to the familiar
kernel formula.  If \(N\subseteq E\) is nuisance and
\(\mathcal A=(A_i)_i\) vanishes on \(N\), then full source recovery is
\(\bigcap_i\ker A_i=N\), while recovery only of a charge \(\chi\) is the
strictly weaker and sharper requirement
\[
                         \ker\mathcal A\subseteq\ker\chi.
\tag{2.0c}\label{eq:charge-kernel}
\]
If the route completes normed spaces, continuous recovery further requires
\(\|\chi(x)\|\le C\|\mathcal A x\|\).  Hilbert lower bounds are therefore
one realization of polarized conservativity, not the universal law.  In a
differential route the factor \(d\) is differential-field reconstruction;
in a jet or Stokes route it separates based formal types; in a cohomological
route it detects exact boundaries.

\subsection{One law, different native factors}

The universal factor \(d\) need not be a bounded operator.  Current routes
ask for the same factorization in different categories:

\begin{center}
\small
\begin{tabularx}{\textwidth}{@{}>{\raggedright\arraybackslash}p{0.18\textwidth}
  >{\raggedright\arraybackslash}p{0.21\textwidth}
  >{\raggedright\arraybackslash}X
  >{\raggedright\arraybackslash}p{0.22\textwidth}@{}}
\toprule
Route & Target charge \(\chi\) & Polarized factor \(d\) & Status / classified boundary \\
\midrule
Calder\'on--Witt Green & absolute orientation and relative incidence escape
  & trace retraction plus Green transgression & exists algebraically; no third boundary-blind channel \\[6pt]

Infinite jets/Stokes & based normal orientation
  & source-pointed Stokes reporter & exact quotient exists; Airy first-deviation reporter is complete in its model \\[6pt]

Gamma differential field & unfolded archimedean charge
  & differential reconstruction before cyclic norm & exponent one is hypertranscendental; even folding changes the differential regime \\[6pt]

Hedenmalm/Bohr & reflected phase or divisor orientation
  & boundary-sensitive transverse reporter & regular phase recovery exists; dense fibres force any separated reporter to be singular \\[6pt]

Hilbert graph & continuous horizontal charge
  & bounded factor on the attained range & exists exactly with a quotient lower bound; two-sided failure creates a bi-dense hyperplane \\[6pt]

Prime Wold--Rees & reflected valuation charge
  & nonlocal reflection sewing & every informative exact factor is nonproduct and beyond all finite Rees orders \\[6pt]

Sonin--Calkin cyclic & commutator orientation
  & graded or noncyclic factor & every cyclic scalar kills the charge; a two-dimensional graded trace detects it \\
\bottomrule
\end{tabularx}
\end{center}

The table is not a list of independent conjectures.  Each row supplies an
observation \(\pi\), a target \(\chi\), and a categorical meaning of
\(\chi=d\pi\).  Positive results construct \(d\) on a restricted carrier;
no-go theorems exhibit a kernel-pair collision or prove that every such
factor is singular.  This is the precise sense in which the route-specific
missing objects are collateral forms of one law.
The four-unit Ward route occupies the same row conceptually: its target is a
signed defect index and its requested factor is a positive,
orientation-sensitive contraction.  Current scalar positivity does not
construct that factor; any completed realization must also satisfy the
lower-bound law.

We now work over an arbitrary field \(k\).  Let \(A_1,A_2,O\) be
\(k\)-vector spaces and suppose
\[
 D:A_1\to A_2,\qquad T:A_2\to O,\qquad
 \jmath:O\to A_2,\qquad H:A_2\to A_1
\tag{2.1}\label{eq:green-data}
\]
are linear and satisfy
\[
 DH+\jmath T=\id_{A_2},\qquad TD=0,
 \qquad T\jmath=\id_O.
\tag{2.2}\label{eq:green-identities}
\]
We call \(O\) the absolute orientation, \(D(A_1)\) the boundary channel,
and \(H\) the chosen Green lift.

For an observer \(F:A_2\to W\), define
\[
 F_{\mathrm{abs}}=(F\jmath)T,
 \qquad
 F_{\mathrm{rel}}=(FD)H.
\tag{2.3}\label{eq:observer-parts}
\]

In the completed setting let \(G,H_1,H_2\) be complex Hilbert spaces and
let \(P:G\to H_1\), \(Q:G\to H_2\) be bounded.  We say that \(P\) is
\emph{uniformly separating} if there is \(C>0\) such that
\[
                         \|x\|\le C\|Px\|
                         \qquad(x\in G).
\tag{2.4}\label{eq:lower-bound}
\]
This is the lower-bound form appropriate to graph realizations.

\begin{theorem}[Green, completion, and grading realizations]
\label{thm:boundary-completeness}
The following statements hold.

\begin{enumerate}[label=\textup{(\alph*)},leftmargin=2.2em]
\item Every linear observer has the exact decomposition
      \[
                         F=F_{\mathrm{abs}}+F_{\mathrm{rel}}.
      \tag{2.5}\label{eq:canonical-split}
      \]
      The absolute part kills every boundary, the relative part kills every
      orientation representative, and these two properties make the split
      unique.

\item The following are equivalent:
      \[
      F_{\mathrm{rel}}=0;
      \qquad FD=0;
      \qquad
      \exists!\,\phi:O\to W\ \text{with }F=\phi T.
      \tag{2.6}\label{eq:factor-classification}
      \]
      Thus \(O\) is the complete absolute invariant; every additional
      observation is precisely boundary sensitivity.

\item Suppose \(P,Q\) are injective and have dense range.  There exists a
      proper closed subspace \(M\subsetneq G\) satisfying
      \[
                   \overline{P(M)}=H_1,
                   \qquad \overline{Q(M)}=H_2
      \tag{2.7}\label{eq:bidense}
      \]
      if and only if neither \(P\) nor \(Q\) is uniformly separating.  In
      the singular case one may take \(M=v^\perp\) for one nonzero vector
      \(v\in G\).

\item Let \(R\) be a ring, let \(e^2=e\), and suppose \(u^2=1\) and
      \(ueu=1-e\).  The orientation \(J=2e-1\) satisfies
      \[
                         J=[eu,ue].
      \tag{2.8}\label{eq:orientation-commutator}
      \]
      Consequently every additive cyclic trace annihilates \(J\).  This
      loss is sharp: for
      \[
      J_0=\begin{pmatrix}1&0\\0&-1\end{pmatrix},\qquad
      U_0=\begin{pmatrix}0&1\\1&0\end{pmatrix},
      \tag{2.9}\label{eq:matrix-pair}
      \]
      one has \(J_0U_0=-U_0J_0\), while the graded trace
      \(A\mapsto\Tr(J_0A)\) takes the nonzero value \(2\) at \(J_0\).
\end{enumerate}
\end{theorem}

Parts (a)--(b) classify source information, part (c) classifies information
lost by singular completion, and part (d) classifies information lost by
scalar cyclicization.  Their conjunction is the promised law: resolve the
relative boundary channel first; retain it by a uniformly separating joint
test; and read its orientation by a graded rather than cyclic scalar.

\begin{corollary}[Exact repair criterion]
\label{cor:repair}
Let \(N\subseteq G\) be a declared nuisance subspace and let
\(A:G/N\to\cH\) be the joint completed analysis map of all retained tests.
The tests recover the completed resolved source stably if and only if
\[
                  \|A[x]\|\ge c\,\|[x]\|
                  \qquad([x]\in G/N)
\tag{2.10}\label{eq:joint-lower-bound}
\]
for some \(c>0\).  Injectivity of \(A\), density of its range, nonflatness
of one test, or completeness of every finite truncation does not imply
\eqref{eq:joint-lower-bound}.
\end{corollary}

\begin{proof}
The first assertion is the definition of stable recovery on the quotient.
For the second, Theorem~\ref{thm:boundary-completeness}(c) supplies an
injective dense-range pair for which a proper closed hyperplane is dense in
every observed completion whenever both legs fail lower bounds.  No finite
or merely algebraic separation statement excludes that hyperplane.
\end{proof}

\section{Absolute orientation and relative Green transgression}
\label{sec:green}

This section proves parts (a)--(b) without topology, positivity, or an inner
product.  The identities are elementary; their consequence is not.  They
show that after an exact source contraction there is no unspecified third
regular channel.

\begin{proposition}[Canonical observation split]
\label{prop:green-split}
Under \eqref{eq:green-identities}, every \(F:A_2\to W\) satisfies
\[
                    F=(F\jmath)T+(FD)H.
\tag{3.1}\label{eq:green-observer-split}
\]
The first summand vanishes on \(\im D\), and the second vanishes on
\(\im\jmath\).
\end{proposition}

\begin{proof}
Applying \(F\) to \(DH+\jmath T=\id\) gives
\[
                     F=(FD)H+(F\jmath)T,
\]
which is \eqref{eq:green-observer-split}.  Since \(TD=0\),
\((F\jmath)TD=0\).  Applying the Green identity to \(\jmath o\), and using
\(T\jmath o=o\), gives \(DH\jmath o=0\); hence \((FD)H\jmath o=0\).
\end{proof}

\begin{proposition}[Uniqueness and factorization]
\label{prop:green-unique}
Suppose \(F=P+S\), where \(PD=0\) and \(S\jmath=0\).  Then
\[
                 P=(F\jmath)T,\qquad S=(FD)H.
\tag{3.2}\label{eq:unique-parts}
\]
Moreover,
\[
 FD=0
 \quad\Longleftrightarrow\quad
 \exists!\,\phi:O\to W\text{ such that }F=\phi T.
\tag{3.3}\label{eq:trace-universal}
\]
\end{proposition}

\begin{proof}
For \(z\in A_2\), apply \(P\) to
\(z=DH z+\jmath Tz\).  The boundary term vanishes, so
\(Pz=P\jmath Tz\).  But \(F\jmath=P\jmath\), because \(S\jmath=0\).
This proves the first identity.  Applying \(S\) to the same decomposition
and using \(PD=0\) proves the second.

If \(FD=0\), Proposition~\ref{prop:green-split} gives
\(F=(F\jmath)T\).  If also \(F=\psi T\), then
\(\psi=\psi T\jmath=F\jmath\), proving uniqueness.  Conversely,
\(F=\phi T\) implies \(FD=\phi TD=0\).
\end{proof}

\begin{remark}
The word ``relative'' is exact.  For an admitted primitive sector
\(\cS\subseteq A_1\), the obstruction to solving \(Db=x\) with
\(b\in\cS\) is the class of \(Hx\) in
\(A_1/(\cS+\ker D)\).  If \(D\) is injective, that class vanishes exactly
when \(Hx\in\cS\).  Thus relative information is canonical Green-primitive
escape, not an additional absolute cohomology class.
\end{remark}

\section{The sharp singular-completion boundary}
\label{sec:singular}

Algebraic separation does not survive completion automatically.  The exact
boundary is most transparent for two Hilbert legs.

\begin{lemma}[Dense image of a hyperplane]
\label{lem:dense-hyperplane}
Let \(T:G\to H\) be bounded with dense range, and let \(v\in G\).  If
\(v\notin\im T^*\), then \(T(v^\perp)\) is dense in \(H\).
\end{lemma}

\begin{proof}
Suppose \(y\perp T(v^\perp)\).  Then \(T^*y\in(v^\perp)^\perp=\C v\),
so \(T^*y=av\).  If \(a\ne0\), then
\(v=T^*(a^{-1}y)\), contrary to the hypothesis.  Hence \(T^*y=0\).
Dense range of \(T\) makes \(T^*\) injective, so \(y=0\).  The orthogonal
complement of \(T(v^\perp)\) is therefore zero, which is equivalent to
density.
\end{proof}

\begin{lemma}[Adjoint range criterion]
\label{lem:adjoint-lower}
Let \(T:G\to H\) be bounded, injective, and have dense range.  Then
\(T\) is uniformly separating if and only if \(T^*\) is onto.
\end{lemma}

\begin{proof}
If \(T^*\) is onto, the open mapping theorem gives a constant \(C\) such
that every \(z\in G\) is \(T^*u\) for some \(u\) with
\(\|u\|\le C\|z\|\).  Then
\[
 \|z\|^2=\langle T^*u,z\rangle
          =\langle u,Tz\rangle
          \le C\|z\|\,\|Tz\|,
\]
which gives \eqref{eq:lower-bound}.  Conversely, a lower bound makes the
range of \(T\) closed.  Dense range then makes \(T\) onto, and the bounded
inverse identifies \((T^{-1})^*\) as a right inverse of \(T^*\).
\end{proof}

\begin{proof}[Proof of Theorem~\ref{thm:boundary-completeness}(c)]
Assume first that neither \(P\) nor \(Q\) is uniformly separating.  By
Lemma~\ref{lem:adjoint-lower}, both \(\im P^*\) and \(\im Q^*\) are proper
linear subspaces of \(G\).  The union of two proper subspaces cannot cover a
vector space, so choose
\[
                  0\ne v\notin\im P^*\cup\im Q^*.
\]
Lemma~\ref{lem:dense-hyperplane} applied to each leg gives
\(\overline{P(v^\perp)}=H_1\) and
\(\overline{Q(v^\perp)}=H_2\).  The hyperplane \(v^\perp\) is closed and
proper.

Conversely, suppose a proper closed \(M\subset G\) has both images dense.
If, for example, \(P\) were uniformly separating, then
\(P|_M\) would also be uniformly separating.  Its range would be closed;
being dense, it would be all of \(H_1\).  For any \(x\in G\), surjectivity
would give \(m\in M\) with \(Pm=Px\), and injectivity would give \(m=x\).
Thus \(M=G\), a contradiction.  The same argument applies to \(Q\).
\end{proof}

\begin{corollary}[Maximal verticality of an over-retained relation]
\label{cor:maximal-verticality}
Let \(Y=\C^{\N}\) carry the product topology and let
\(\mathcal R\subseteq X\times Y\) be a closed linear relation.  If its
multivalued part projects onto every leading finite coordinate block, then
\[
                 \operatorname{mul}\mathcal R=Y,
                 \qquad
                 \mathcal R=\operatorname{dom}\mathcal R\times Y.
\tag{4.1}\label{eq:maximal-vertical}
\]
Thus the retained relation carries no output information.
\end{corollary}

\begin{proof}
Surjectivity onto every leading finite block makes
\(\operatorname{mul}\mathcal R\) dense in the product topology.  The
multivalued part is the intersection of the closed relation with
\(\{0\}\times Y\), hence is closed.  It is therefore all of \(Y\).  If
\((x,y_0)\in\mathcal R\) and \(y\in Y\), then
\((0,y-y_0)\in\mathcal R\), so \((x,y)\in\mathcal R\).
\end{proof}

The corollary is the infinite-coordinate extreme of the same law.  A
relation that retains all finite outputs without a separating global bound
need not be ``highly multivalued''; it can be maximally vertical.

\subsection{The complete prime-normal source: nonlocal and singular}
\label{subsec:nonlocal-singular}

The Riemann source makes the abstract boundary quantitative.  Put
\[
 H_m(s)=\frac1{\zeta((m+1)s)}-\frac1{\zeta(ms)}
\]
and, for \(K\ge1\) and \(\Re s>0\), define the complete shell vector in \(\ell^1(\N)\) by
\[
 X_K(s)=\bigl(\zeta(Ks)^{-1},H_K(s),H_{K+1}(s),\ldots\bigr).
\tag{4.2}\label{eq:complete-source-vector}
\]
The coordinates telescope to one.  At each finite \(K\), the map
\(s\mapsto X_K(s)\) is injective; for \(K\ge12\) it contains a nonzero
coordinate carrying the complete \(\Xi\)-divisor order at every point with
real part at least one half.  Cofinally, however,
\[
                    X_K(s)\longrightarrow (1,0,0,\ldots)
                    \qquad(K\to\infty).
\tag{4.3}\label{eq:cofinal-collapse}
\]

For the reflected valuation sectors, let
\[
 V_N=\prod_{p<N}\left(1-\frac1p\right)
\tag{4.4}\label{eq:wold-volume}
\]
and, for \(|c|<\tfrac12\), let their squared finite-prime Hellinger affinity
be
\[
 B_c(N)=\prod_{p<N}
 \frac{(1-p^{-1})(1-p^{-(1-2c)})}
      {(1-p^{-(1-c)})^2}.
\tag{4.5}\label{eq:hellinger-sector}
\]

\begin{theorem}[Prime-normal nonlocal--singular law]
\label{thm:nonlocal-singular}
The complete source satisfies all of the following.
\begin{enumerate}[label=\textup{(\roman*)},leftmargin=2.2em]
\item For distinct \(s,t\) with positive real part,
      \[
        \|X_K(s)-X_K(t)\|^{-1}\longrightarrow\infty.
      \tag{4.6}\label{eq:inverse-conditioning}
      \]
      Consequently no family of observers with one uniform Lipschitz
      constant can retain a fixed positive separation between \(s\) and
      \(t\) at all cofinal depths.
\item A closed linear relation that retains every countable prime-normal
      packet in the product topology is maximally vertical.  Any informative
      closed relation carrying the same packets must use a topology strictly
      finer than the product topology; its information is invisible to every
      finite set of prime coordinates.
\item Let \(0<c<\tfrac12\), and let
      \(S_N:\R\to\R\) be bounded linear maps satisfying
      \(S_N(B_c(N))=B_{-c}(N)\).  Then, for every \(r\in\N\),
      \[
                  V_N^{-r}=o(\|S_N\|).
      \tag{4.7}\label{eq:beyond-rees}
      \]
      Thus exact reflection sewing grows beyond every finite Wold--Rees
      order.
\end{enumerate}
An informative reflected completion is therefore simultaneously nonlocal in
topology and essentially singular in operator order.
\end{theorem}

\begin{proof}
The telescoping identity gives unit source mass and
\eqref{eq:cofinal-collapse}; the head tends to one and the \(\ell^1\)-norm
of the remaining shell tail tends to zero.  Finite-depth injectivity makes
the norm in \eqref{eq:inverse-conditioning} positive, so convergence of the
difference to zero implies divergence of its reciprocal.  The Lipschitz
assertion follows immediately.

The countable prime-normal packet has surjective projection onto every
finite leading block.  Its span is dense in product topology.  Closedness
and Corollary~\ref{cor:maximal-verticality} then force maximal verticality.
Hence an informative topology retaining the packets must be strictly finer.

For \(c>0\), \(B_c(N)\) is flat beyond every power of \(V_N\), whereas
\(B_{-c}(N)\) has first Wold order.  The operator inequality
\[
              B_{-c}(N)=|S_N(B_c(N))|
                 \le \|S_N\|B_c(N)
\]
therefore gives \eqref{eq:beyond-rees} at every order.  Combining the last
two conclusions proves the final statement.
\end{proof}

The theorem is strong but route-specific.  It constrains every approach
that completes the countable prime-normal source and exactly sews its
reflected valuation sectors.  It does not rule out a differential,
cohomological, or Stokes argument that obtains orientation before entering
that completion.

\section{Why scalar traces lose orientation}
\label{sec:orientation}

The completion theorem explains hidden sectors but not the parity of the
missing information.  That parity is fixed by a small crossed-product
calculation.

\begin{proposition}[Commutator orientation]
\label{prop:commutator-orientation}
Let \(R\) be a ring.  Suppose \(e^2=e\), \(u^2=1\), and \(ueu=1-e\).
Put
\[
                         X=eu,\qquad Y=ue,
                         \qquad J=2e-1.
\]
Then \(XY=e\), \(YX=1-e\), and \(J=XY-YX\).  If
\(\tau:R\to B\) is additive and cyclic, meaning
\(\tau(ab)=\tau(ba)\), then \(\tau(J)=0\).
\end{proposition}

\begin{proof}
Using the three relations,
\[
 XY=euu e=e^2=e,
 \qquad
 YX=ueeu=ueu=1-e.
\]
Thus \(XY-YX=2e-1=J\).  Cyclicity gives
\(\tau(J)=\tau(XY)-\tau(YX)=0\).
\end{proof}

The involution exchanges the two complementary sectors, while \(J\)
records their signed difference.  Proposition~\ref{prop:commutator-orientation}
shows that cyclic scalarization is structurally incapable of reading that
difference.  Positivity does not alter the conclusion: a positive cyclic
trace is still cyclic.

\begin{proposition}[Minimal graded recovery]
\label{prop:graded-recovery}
For the matrices in \eqref{eq:matrix-pair},
\[
 J_0^2=U_0^2=I,\qquad J_0U_0=-U_0J_0,
 \qquad \Tr(J_0J_0)=2.
\tag{5.1}\label{eq:graded-detection}
\]
Thus two dimensions already suffice to recover an orientation invisible to
ordinary trace.
\end{proposition}

\begin{proof}
Direct matrix multiplication gives the first three identities.  Since
\(J_0^2=I\), the last value is \(\Tr(I)=2\).
\end{proof}

\begin{corollary}[Necessary form of an orientation repair]
\label{cor:orientation-repair}
Any observer required to distinguish the two sectors exchanged by \(u\)
must violate at least one hypothesis of Proposition~\ref{prop:commutator-orientation}:
it must be noncyclic, graded, boundary-sensitive, or vector-valued before
its final scalar readout.
\end{corollary}

This is a necessity theorem, not an assertion that every noncyclic observer
works.  The completion law adds the second requirement: the joint graded
observer must also satisfy a lower bound on the resolved quotient.

\section{The resolved arithmetic source}
\label{sec:source}

We now identify the source to which the abstract law applies.  The local
arithmetic construction has degree-one space
\[
                  \cA_p^1=\cQ_p\oplus\cW_p,
\]
where \(\cQ_p\) is the Meyer--Calder\'on defect sector and \(\cW_p\) is the
directed rational-Witt packet.  Its alternating cross-effect is a
vector-valued cup
\[
            \smile:\Lambda^2\cA_p^1\longrightarrow\cA_p^2.
\tag{6.1}\label{eq:source-cup}
\]
The cup is onto, its vector-valued radical is zero, and its kernel
\(\cI_p\) is the exact incidence relation space.  There is an explicit
section \(\sigma_p\) and Green retraction \(h_p\), yielding a split exact
complex
\[
 \cI_p\longrightarrow\Lambda^2\cA_p^1
        \xrightarrow{\smile}\cA_p^2.
\tag{6.2}\label{eq:local-source-complex}
\]
The same construction holds on the all-prime finite-support direct sum.

A normalized Haar coefficient \(\tau:\cA^2\to\C\) adds the absolute
orientation line.  On \(\cA^2\times\C\), put
\[
 D w=(\smile w,-\tau(\smile w)),
 \qquad
 T(z,c)=\tau(z)+c,
 \qquad
 \jmath(o)=(0,o).
\tag{6.3}\label{eq:oriented-source}
\]
The Green lift inherited from the section satisfies
\[
                         DH+\jmath T=\id,
                         \qquad TD=0,
                         \qquad T\jmath=\id.
\tag{6.4}\label{eq:source-green}
\]

\begin{theorem}[Arithmetic source completeness]
\label{thm:arithmetic-source}
For the local and all-prime sources:
\begin{enumerate}[label=\textup{(\roman*)},leftmargin=2.2em]
\item the unoriented source is split exact;
\item the oriented source has cohomology \((0,0,\C)\);
\item every top-source observer is uniquely the sum of its orientation
      compression and its boundary Green transgression;
\item an observer insensitive to incidence boundaries factors uniquely
      through the one-dimensional trace \(T\);
\item after restricting incidence primitives to a sector
      \(\cS\subseteq\cI\), the remaining class vanishes exactly when the
      canonical Green primitive belongs to \(\cS\).
\end{enumerate}
\end{theorem}

\begin{proof}
Split exactness follows from the cup section and Green retraction.  The
oriented boundary is the graph of the negative Haar coefficient, so its
cokernel is measured exactly by \(T\), giving the single top class.
Equations \eqref{eq:source-green} instantiate
Propositions~\ref{prop:green-split} and~\ref{prop:green-unique}.  The
relative assertion follows by applying the Green contraction to a cocycle:
its primitive class in \(\cI/(\cS+\ker D)\) vanishes exactly when it has a
representative in \(\cS\); here the incidence inclusion is injective.
\end{proof}

The theorem rules out a frequent misdiagnosis.  The source is not missing a
regular scalar invariant beyond orientation.  Such an invariant cannot
exist: it would have to detect an exact boundary and is therefore a relative
transgression.  The unresolved information is not absent from the source;
it is lost when relative transgression is scalarized or completed without a
lower bound.

\subsection{Infinite jets: position and based velocity}

The source split is not confined to a finite combinatorial model.  Let
\(G\) be a commutative group scheme over a field \(K\), and let \(R\) be a
\(K\)-algebra.  Write \(J^nG(R)\) for its truncated algebraic jets and
\(J^\infty G(R)\) for coherent families under all truncation maps.
Evaluation at the base point has a constant section.

\begin{theorem}[Infinite Maurer--Cartan splitting]
\label{thm:infinite-jet-split}
There are canonical group equivalences
\[
 J^\infty G(R)
   \simeq G(R)\times J^\infty_eG(R),
\tag{6.5}\label{eq:infinite-split}
\]
\[
 J^\infty G(R)/G(R)
   \simeq J^\infty_eG(R),
\tag{6.6}\label{eq:constant-quotient}
\]
and, for every subgroup \(L\le G(R)\),
\[
 J^\infty G(R)/L
   \simeq (G(R)/L)\times J^\infty_eG(R).
\tag{6.7}\label{eq:subgroup-quotient}
\]
Moreover, every formally \(\acute{e}\)tale affine morphism preserves all
finite truncated arcs by pullback along the underlying point.
\end{theorem}

\begin{proof}
Evaluation and constant-jet insertion form a split homomorphism.  In a
commutative group, every split homomorphism is the product of its image and
kernel, giving \eqref{eq:infinite-split}; quotienting the constant image
gives \eqref{eq:constant-quotient}, and quotienting a subgroup of the
constant factor gives \eqref{eq:subgroup-quotient}.  Coherence is preserved
coordinatewise by truncation.  For a formally \(\acute{e}\)tale morphism,
the augmentation ideal of \(R[X]/(X^{n+1})\) is nilpotent.  The defining
infinitesimal lifting property therefore makes the arc pullback map
bijective at every order \(n\).
\end{proof}

Equations \eqref{eq:infinite-split}--\eqref{eq:subgroup-quotient} identify
the correct pre-divisor carrier.  Constant position is a declared nuisance
sector; the full based infinite jet is the resolved Maurer--Cartan datum.
A consumer that squares, traces, or completes before acting on this quotient
can preserve divisor support while losing its normal orientation.

\section{The Gamma boundary is unfolded differential data}
\label{sec:gamma}

The Gamma factor is often treated as a scalar normalization appended after
the arithmetic source.  The canonical-product family shows why that order
is dangerous.

For \(p>0\), let \(W_p\) be the canonical product with zeros
\(-n^p\), \(n\ge1\), using the minimal primary genus.  At exponent one the
product is exactly a reciprocal Gamma gauge.

\begin{theorem}[Gamma and even-fold dichotomy]
\label{thm:gamma-even}
For every \(z\in\C\),
\[
                  W_1(z)=
                  \frac{e^{-\gamma z}}{\Gamma(z+1)},
\tag{7.1}\label{eq:W1-gamma}
\]
where \(\gamma\) is Euler's constant.  The entire function \(W_1\) is
hypertranscendental over \(\C(z)\).  In contrast, for every \(m\ge1\),
\(W_{2m}\) is differentially algebraic.  More precisely, if
\(\alpha\) is a primitive \(2m\)-th root of unity and
\[
                  S(w)=\prod_{n\ge1}
                        \left(1-\frac{w^2}{n^2}\right),
\]
then
\[
 W_{2m}(-w^{2m})=
        \prod_{k=0}^{m-1}S(\alpha^kw).
\tag{7.2}\label{eq:even-sine-factorization}
\]
For every \(q\ge1\), \(W_{1/q}\) is also hypertranscendental, and for a
primitive \(q\)-th root \(\omega\),
\[
 \prod_{j=0}^{q-1}W_{p/q}(\omega^jz)
   =W_p((-1)^{q+1}z^q)
\tag{7.3}\label{eq:cyclic-norm}
\]
for all positive integers \(p,q\).
\end{theorem}

\begin{proof}
The genus-one Weierstrass factors give
\[
 W_1(z)=\prod_{n\ge1}
       \left(1+\frac zn\right)e^{-z/n}.
\]
Multiplying the finite products by the standard Gamma approximants and
passing to the limit yields \eqref{eq:W1-gamma}, including the negative
integers by continuity.  If \(W_1\) were differentially algebraic, then
multiplication by the differential-algebraic exponential gauge
\(e^{\gamma z}\) would make \(1/\Gamma(z+1)\) differentially algebraic,
contradicting the differential transcendence of the Gamma orbit.

For \(p=2m\), the primary genus is zero.  Factor
\[
 1-\frac{w^{2m}}{n^{2m}}
   =\prod_{k=0}^{m-1}
      \left(1-\frac{(\alpha^kw)^2}{n^2}\right)
\]
and pass to convergent products to obtain
\eqref{eq:even-sine-factorization}.  Each sine product satisfies an
algebraic differential equation, and differential-algebraic functions are
closed under finite products and power substitution; hence \(W_{2m}\) is
differentially algebraic.

For \eqref{eq:cyclic-norm}, multiply the \(q\) root-of-unity packets at each
zero index.  The exponential parts of the primary factors cancel by the
root-of-unity sums, leaving exactly the primary factor of exponent \(p\) at
\((-1)^{q+1}z^q\).  Convergent product uniqueness gives the identity.
Differential algebraicity descends through finite power substitution; its
contrapositive transfers hypertranscendence of \(W_1\) to every
\(W_{1/q}\).
\end{proof}

The theorem is a control experiment for the polarized conservativity law.
Canonical-product form, negative real zero support, and finite cyclic
factorization do not determine differential complexity.  The unfolded
Gamma sector can be hypertranscendental while an even fold is a finite
product of differential-algebraic sine sectors.  Thus Gamma data should be
tested before a parity fold or scalar cyclic norm.  This is the analytic
counterpart of Proposition~\ref{prop:commutator-orientation}: scalar
aggregation preserves a divisor pattern but can erase the structure needed
to orient it.

\section{What regular tomography already recovers}
\label{sec:regular}

The missing Riemann datum cannot be blamed on a shortage of regular tests.
Two exact completeness results make this precise.

First, let
\[
 f_\lambda(z)=\sum_{j\in I}a_j\lambda_j e^{i\omega_jz},
 \qquad |\lambda_j|=1,
\tag{8.1}\label{eq:bohr-carrier}
\]
be normally summable on compact sets, with nonzero amplitudes and distinct
real frequencies.  The map from the phase torus \(\mathbb T^I\) to the
compact-open space of entire functions is a closed embedding.  If all
frequencies are positive and \(c\ne0\), equality of any fixed nonzero powers
\[
                    (c+f_\lambda)^n=(c+f_\mu)^n
\]
still forces \(\lambda=\mu\), provided the denominator carrier is
zero-free.  The asymptotic value at the upper imaginary end fixes the
otherwise possible root of unity.  Regular phase data therefore survive
ordinary power observation when an oriented base point is retained.

Second, let \(U_j(t)\) be bounded pulse operators on a normed space, with a
common invariant domain \(D\) and possibly unbounded generators \(A_j\).
Assume only
\[
 \left.\frac d{dt}U_j(t)v\right|_{t=0}=iA_jv
 \qquad(v\in D).
\tag{8.2}\label{eq:pulse-derivative}
\]
For a continuous linear functional \(\omega\), one derivative in each
coupling recovers the complete ordered correlation:
\[
 \left.\partial_{t_0}\cdots\partial_{t_{n-1}}
   \omega\!\left(U_0(t_0)\cdots U_{n-1}(t_{n-1})v\right)
 \right|_{t=0}
   =i^n\omega(A_0\cdots A_{n-1}v).
\tag{8.3}\label{eq:ordered-tomography}
\]
There are explicit pulse families satisfying
\eqref{eq:pulse-derivative} whose diagonal generator is genuinely
unbounded.  Hence the statement is not a bounded-operator surrogate.

Equations \eqref{eq:bohr-carrier}--\eqref{eq:ordered-tomography} delimit the
frontier.  Regular commutative phases and regular noncommutative ordered
moments are already fully observable.  What is absent is not another
regular moment family but an orientation-complete boundary test.  If a route
passes through the completion of Theorem~\ref{thm:boundary-completeness}(c),
that test must also be uniform and singular; other routes retain their native
notion of completeness.

\section{The Riemann frontier}
\label{sec:riemann}

Let \(\Xi\) denote the completed Riemann function and let \(\rho\) be a
nontrivial zero.  Its horizontal charge is
\[
                         a_\rho=\Re\rho-\frac12.
\tag{9.1}\label{eq:horizontal-charge}
\]
Functional-equation reflection changes its sign.  The logarithmic Green
connection
\[
                    \cG_\Xi(s)=-2\frac{\Xi'(s)}{\Xi(s)}
\]
has, after multiplication by the balanced Euler field, residue
\[
 \operatorname*{res}_{s=\rho}
 \left(\frac{s-\frac12}{2}\cG_\Xi(s)\right)
       =-m_\rho\!\left(\rho-\frac12\right).
\tag{9.2}\label{eq:outer-charge}
\]
Thus the unsquared Green connection already carries the full complex outer
charge.  Squaring the coordinate preserves the divisor location but erases
the sign of \(a_\rho\).

The arithmetic source of Section~\ref{sec:source} supplies exact incidence,
one absolute orientation line, and the relative Green primitive.  The jet
split supplies the based infinite-order carrier.  Section~\ref{sec:regular}
shows that regular phase and ordered-moment tomography are complete.
Theorems~\ref{thm:resolved-test}, \ref{thm:boundary-completeness}, and
\ref{thm:gamma-even} then give the following common specification.

\begin{definition}[Admissible Riemann boundary test]
\label{def:admissible-test}
An admissible test \(\Sigma\) is a source-derived observation on the resolved
quotient in the native category of a proposed route such that:
\begin{enumerate}[label=\textup{(\roman*)},leftmargin=2.2em]
\item it is defined before cyclic trace, parity folding, or divisor
      scalarization;
\item it is orientation-sensitive---reflection odd and graded, or
      equivalently noncyclic/vector-valued before final readout---so it does
      not annihilate the commutator orientation;
\item together with the existing source tests, it separates the quotient by
      constant position and exact incidence; if the route passes to a normed
      completion, the joint analysis is additionally bounded below;
\item on a divisor germ at \(\rho\), its graded boundary value is a nonzero
      fixed multiple of \(m_\rho a_\rho\);
\item the source Green identity forces that graded boundary value to vanish.
\end{enumerate}
\end{definition}

\par\medskip

\begin{theorem}[Riemann closure theorem]
\label{thm:riemann-closure}
Existence of an admissible Riemann boundary test implies the Riemann
hypothesis.
\end{theorem}

\begin{proof}
Let \(\rho\) be any nontrivial zero.  Property (v) makes the graded boundary
value zero.  Property (iv), with \(m_\rho>0\), gives
\(a_\rho=0\).  Hence \(\Re\rho=1/2\).  Since \(\rho\) was arbitrary, every
nontrivial zero lies on the critical line.
\end{proof}

The theorem is deliberately not advertised as a proof of the Riemann
hypothesis: Definition~\ref{def:admissible-test} includes the producer
property that remains to be constructed.  Its content is the minimality of
the specification.  Removing any item returns a failure already classified
above:

\begin{center}
\begin{tabularx}{\textwidth}{@{}>{\raggedright\arraybackslash}p{0.25\textwidth}
  >{\raggedright\arraybackslash}X
  >{\raggedright\arraybackslash}p{0.29\textwidth}@{}}
\toprule
Omitted requirement & What survives & Exact obstruction \\
\midrule
Resolved source quotient & constant position or exact incidence & nuisance can mimic divisor data \\
Graded, reflection-odd readout & cyclic scalar trace & commutator orientation maps to zero \\
Completed lower bound & injective dense regular legs & proper closed bi-dense hyperplane \\
Pre-fold Gamma action & scalar norm or even folding & unfolded differential data can collapse \\
Source-forced vanishing & divisor equivalence only & no zero exclusion \\
\bottomrule
\end{tabularx}
\end{center}

This table is the practical conclusion.  Hedenmalm dense radicals,
Wold--Poisson ancestry, Sonin/Calkin consumers, finite Rees moderation, and
four-unit Ward positivity are not separate mysteries.  They are shadows of
one of three losses: the nuisance quotient is unresolved, the native test
family is incomplete (or not uniformly separating after completion), or its
scalarization kills the relative orientation.  Their
positive survivors identify pieces of Definition~\ref{def:admissible-test},
but none supplies all five properties.

\subsection{Why the result changes the search}

The law removes several broad search spaces.

\begin{enumerate}[leftmargin=2.1em]
\item \emph{Additional regular sources.}  The arithmetic source already has
      exact incidence, one absolute orientation, and complete relative
      transgression.  Infinite jets add based velocity without creating an
      unexplained regular class.
\item \emph{Additional scalar positivity.}  A cyclic scalar cannot see the
      grading \(2e-1\).  Positivity of such a scalar cannot force its sign.
\item \emph{Bounded completion tricks.}  Two singular dense legs admit a
      codimension-one closed sector invisible in both completions.  A finite
      Rees-moderate or bounded graded operator cannot repair a flat radical
      without a new lower bound.
\item \emph{Divisor-preserving folds.}  Gamma canonical products show that
      finite folding can preserve an exact zero pattern while changing the
      differential-algebraic regime.  Divisor preservation is not resolved
      information preservation.
\end{enumerate}

What remains is not restricted to an operator-theoretic construction.
Every route must supply an orientation-complete source test satisfying the
five checkable properties in Definition~\ref{def:admissible-test}.  A Stokes
route may do so by a first-deviation sign law; a differential route by an
unfolded Gamma invariant; a cohomological route by a nonzero relative
transgression.  Only a route that enters completed prime-normal geometry is
forced to realize the test as a nonlocal singular graph with a uniform
quotient lower bound.

\section{Conclusion}

The polarized conservativity law says that information has three lives.
Algebraically, it is an absolute orientation plus a relative Green
transgression.  Topologically, it survives completion exactly through a
uniformly separating joint test.  Parity-sensitive information survives
scalarization only through a graded, noncyclic readout.

For the Riemann divisor, current source constructions already solve the
first life.  Regular Bohr and ordered-pulse tomography solve ordinary
separation.  The Gamma canonical-product dichotomy explains why folding the
archimedean factor is irreversible, and the commutator theorem explains why
cyclic positivity cannot restore its orientation.  The two-sided singular
criterion then locates the precise boundary between formal reconstruction
and completed information.

The resulting research target is sharply specified without choosing a
Hilbert--P\'olya route: act on the based infinite Green primitive before
scalarization, retain the Gamma/co-Poisson reflection grade, and construct a
complete native test that turns the outer residue
\(-m_\rho(\rho-\tfrac12)\) into a source-forced vanishing boundary value.
In a completed route ``complete'' means a joint lower bound and entails the
nonlocal--singular costs proved above; in algebraic, differential, and Stokes
routes it has their native separation meaning.  Any such construction would
prove the Riemann hypothesis.  The present paper classifies the common
requirement and the exact cost paid by each current realization.

\end{document}
